\section{Experimental Protocol}
\label{sec:experiments}

\paragraph{Metrics.}
We report MSE on log-RV per horizon (literature comparability); \Rtwo{} vs.\
persistence (\citealp{gukellyxiu2020};
$1 - \sum(y-\hat y)^2 / \sum(y - \hat y_{\mathrm{persist}})^2$), the headline metric,
which is not gameable by identity memorisation the way raw MSE is; MAE; and
cross-sectional Spearman rank correlation per calendar quarter. All metrics are
reported on all three targets ($v$, \dv, HAR-residual) $\times$ all splits, and
carry the horizon convention as a column.

\paragraph{Statistical rigor.}
Every learned model is run with $\geq$5 seeds (mean $\pm$ std). We use
Diebold--Mariano tests \citep{dm1995} (HLN small-sample correction) on per-call
squared-error differentials vs.\ persistence and vs.\ the best preceding stage;
cluster-bootstrap confidence intervals (1{,}000 resamples) clustered by ticker and
by calendar quarter; and Holm-corrected $p$-values across horizons for confirmatory
claims. Confirmatory vs.\ exploratory analyses are labelled in advance: confirmatory
analyses are Stage-$k$ vs.\ Stage-1 DM tests on \dv{} at each $\tau$, Stage~4 vs.\
best unimodal, the temporal-vs-\tickerdisjoint{} gap per stage, and---for each
reproduced prior model---the published-split vs.\ leakage-proof-split gap; all
per-year/regime, COVID-exclusion, prepared-vs-Q\&A, call-length, and
gender-confound analyses are exploratory.

\paragraph{Identity-control suite.}
\label{sec:controls-protocol}
The signature diagnostic of this benchmark probes whether content models read
content or identity, via four controls. \emph{(1) Ticker-only model}: predict the
per-ticker target mean (global mean for unseen tickers); if a content model matches
it on level-$v$, that model reads identity. \emph{(2) Same-ticker transcript
shuffle}: replace each test call's transcript with a different call from the
\emph{same} ticker; if performance is unchanged the model reads identity, not
content, whereas a \emph{global} shuffle (any other call) isolates pure overfitting.
\emph{(3) Identity probe}: a linear probe predicting the ticker from frozen call
embeddings; high accuracy demonstrates identity is encoded. \emph{(4)
Ticker-disjoint split} results reported beside temporal-split results in every
table. The suite runs unchanged on the reproduced prior models
(\S\ref{sec:repro-protocol}).

\paragraph{Post-cutoff experiments.}
Earnings25 (\S\ref{sec:corpora}) supports two experiments that the 2019--2021
corpora cannot. \emph{Clean audio}: the frozen Stage-3 audio ladder (same
extractors, same heads, same controls) is re-run on source-quality 2025 recordings,
to test whether the audio-inert finding of \S\ref{sec:results} is a property of the
task or of FinCall-Surprise's uniform 40\,kbps transcodes. \emph{Lookahead}: the
entire frozen pipeline---no retraining, no post-hoc threshold changes (rule fixed
in advance)---is evaluated on calls held after the knowledge cutoff of every model
used; disproportionate degradation relative to the 2019--2021 results would be
contamination evidence (\S\ref{sec:results-clean-audio}, \S\ref{sec:results-lookahead}).

\paragraph{Reproducibility.}
All result tables regenerate byte-identically from committed CSV artifacts via
\texttt{ecvol report}, guarded by a CI test. Every result CSV carries a committed
run manifest---the resolved command config and its hash, git SHA, seeds,
environment fingerprint, and the file's SHA-256---and both \texttt{ecvol report}
and CI refuse a result file whose bytes match no manifest; one committed config
per result-producing command is the record of how it was run. Every data file
carries a SHA-256 manifest entry; targets are unit-tested against hand-computed
values; and the splits are guarded by leakage assertions in CI. Every result that
predated this mechanism has since been re-run by its producing command and
reproduced its committed file byte for byte, five-seed MLP heads included.
\texttt{REPRODUCE.md} lists the one command per
table; a clean-machine regeneration of Table~\ref{tab:baselines} is part of the
release checklist.
